\section{Finite projected Heisenberg signatures under a whole-source floor}
\label{sec:heisenberg-finite-signature}

\paragraph{Status and provenance.}
D-TD63, question Q2fs, is separate from D-TD61.
The parent proposed the three-entry and four-entry quantitative
continuations of the orbit-span bound.
The author checked those inequalities, supplied the complete
low-exponent combination and the faithful cyclic-module order bound.
The complete author proof, code and first-output readings pass.
Complete root and independent whole-proof reviews also pass.
The conclusion bounds selected projective inertia orders, not the
Heisenberg exponent or the original scalar periods.

\paragraph{The same actual sphere source.}
Let a connected finite \(J\)-Galois cover of \(\mathbf P^1\) have a
faithful irreducible complex coefficient representation \(V\), of
dimension \(e\), and actual generating product-one inertia
\(T_1,\ldots,T_r\).  Suppose an integer \(M\ge1\) satisfies
\[
 T_i\ne I,\qquad r\ge12,\qquad
 \sum_{i=1}^r\operatorname{rank}(T_i-I)=2e+6,\qquad
 \frac{\operatorname{rank}(T_i-I)}e\ge\frac1M
       \quad(1\le i\le r).
 \tag{HFS.1}
\]
A real-valued floor parameter can first be replaced by its ceiling.
The condition is on every original genuine entry; no such condition
is inferred from a selected factor alone.

Suppose \(V=W\otimes U\) is preserved factorwise, where the actual
finite projective image \(P\) on \(W\) is primitive and contains a
normal, nondegenerate projective Heisenberg subgroup
\[
 Q\simeq\mathbf F_p^{\,2m},\qquad m\ge1,\qquad
                 a=\dim W=p^m.
 \tag{HFS.2}
\]
The represented action of \(Q\) on \(W\) is irreducible.
The entire complementary factor \(U\) is arbitrary.
Let \(H\subset\operatorname{Sp}(Q)\) be the actual label image.
Assume separately that \(Q\) is irreducible as an
\(\mathbf F_pH\)-module.  This last property is available for the
minimal-normal Heisenberg extraction of TD12, but is not inferred
merely from primitivity of an arbitrary normalizer.

Write \(y_i\in P\) for the projective inertia, \(g_i\in H\) for its
label action, and \(D_i=\deg\operatorname{minpoly}_{\mathbf F_p}(g_i)\).
Let \(k\) count the nonidentity \(y_i\), retaining all \(r\) entries
in (HFS.1).  In particular \(g_i\) can be identity for \(y_i\ne1\).
Define the following exact integers:
\[
 \begin{split}
 \ell(M)&=\min\{j\in\mathbf Z_{\ge0}:2^j\ge8M\},\\
 B(M)&=\max\{7743,M\},\\
 C(M)&=\max\{64M^2,\ 318\ell(M),\ 2544\},\\
 O(M)&=B(M)^{\,C(M)^2+1}.
 \end{split}
 \tag{HFS.3}
\]

\begin{theorem}[D-TD63]
In this actual-source setting,
\[
 p\le B(M),\qquad D_i\le C(M),\qquad
       \operatorname{ord}(y_i)<O(M)\quad(1\le i\le r),
       \qquad r\le5M.
 \tag{HFS.4}
\]
Thus the ordered ramification-order signature of this selected
actual projective quotient, with order-one entries retained at
the original branch values, lies in an explicitly finite set.

More precisely, when \(k=3\) one has
\[
                D_i<\max\{64M^2,318\ell(M)\}
                    \quad(1\le i\le r).
 \tag{HFS.5}
\]
When \(k=4\), one has \(D_i\le2544\), and if \(m\ge1272\)
then every \(D_i\le159\).
The four-entry assertion itself needs neither \(r\ge12\) nor the
whole-source \(M\)-floor; it uses the actual \(2e+6\) budget and
the irreducible-label hypotheses of TD61.

No bound on \(m\), on \(|P|\), on the number of possible quotient
groups or Nielsen tuples, or on the full original orders
\(\operatorname{ord}(T_i)\) is asserted.
\end{theorem}

\begin{proof}
\emph{Step 1: quantitative decay from the cyclic-module estimate.}
The hypotheses of TD61 apply with the entire original partner.
For every actual factor operator \(A_i\), its all-target
eigenspace estimate is
\[
 1-\frac{\nu_W(A_i)}a
 \le
 \left(\frac1{R+1}+\frac43p^{\,2R-D_i}\right)^{1/2}
       \quad(0\le R<D_i).
 \tag{HFS.6}
\]
This is the full-original-period bound, including all projectively
identity powers; no label-only period is substituted.

If \(D_i\ge64M^2\), choose \(R=16M^2\).
Since \(p\ge2\) and \(2^{32M^2}>32M^2\), one has
\[
 \frac1{16M^2+1}+\frac43p^{-32M^2}
 <
 \frac1{16M^2}+\frac1{24M^2}
 =\frac5{48M^2}<\frac1{M^2}.
 \tag{HFS.7}
\]
The elementary inequality used here follows, for example, by
the binomial estimate \(2^s=(1+1)^s\ge1+s>s\) at
the positive integer \(s=32M^2\).
Thus \(D_i\ge64M^2\) implies
\(\nu_W(A_i)/a>1-1/M\).

TD61 also gives, independently of \(M\), at most one nonidentity
projected entry with \(D_i\ge160\), and for every pair of
distinct nonidentity projected entries \(i,j\),
\[
 2m\le159\max\left\{1,\
       \sum_{\substack{\ell:y_\ell\ne1\\\ell\ne i,j}}
                    \operatorname{rank}_{\mathbf F_p}(g_\ell-I)
                    \right\}.
 \tag{HFS.8}
\]
The maximum with one is essential at the zero-rank exception.

\emph{Step 2: all three-entry minimal polynomials.}
Suppose \(k=3\) and \(m\ge159\ell(M)\).
Here \(\ell(M)\ge3\), so \(2m/159>1\).
Apply (HFS.8) to the other two indices.  For each of the three
nonidentity projected entries this gives
\[
       s_i:=\operatorname{rank}_{\mathbf F_p}(g_i-I)
                    \ge\frac{2m}{159}.
 \tag{HFS.9}
\]
The TD08 trace-and-support inequality, valid for every nonzero
label rank and every eigenvalue target, gives
\[
 \frac{\nu_W(A_i)}a
       \ge\frac{1-p^{-s_i/2}}2
       \ge\frac12-\frac12\,2^{-m/159}.
 \tag{HFS.10}
\]
If one of these three entries had \(D_i\ge64M^2\), use (HFS.7)
on it and (HFS.10) on the other two.
Since \(2^{-m/159}\le2^{-\ell(M)}\le1/(8M)\),
their total normalized cost would be strictly greater than
\(2-1/M-1/(8M)\).
The other \(r-3\ge9\) original entries are still genuine
coefficient operators.  They retain their actual cost at least
\(1/M\) each, even though their selected projective image is identity.
Convolution with the whole \(U\) therefore gives
\[
 2+\frac6e
   >2-\frac1M-\frac1{8M}+\frac9M
    =2+\frac{63}{8M}.
 \tag{HFS.11}
\]
On the other hand \(e\ge p^m\ge2^m\ge2^{\ell(M)}\ge8M\),
so \(6/e\le3/(4M)<63/(8M)\), a contradiction.
All three \(D_i\) are therefore less than \(64M^2\) in this range.
If \(m<159\ell(M)\), the elementary dimension bound gives
\(D_i\le2m<318\ell(M)\).
An omitted entry has \(g_i=I\) and \(D_i=1\).
This proves (HFS.5) for every original index.

\emph{Step 3: all four-entry minimal polynomials.}
Now suppose \(k=4\) and \(m\ge1272\).
If \(D_i\ge160\) for some entry, its normalized projective support
is strictly greater than \(22/25\), by TD61.
For any two of the other three label ranks, apply (HFS.8)
by excluding \(i\) and the third remaining index.
Since \(2m/159>1\), every such pair has sum at least \(2m/159\).
Consequently at least two of those three ranks are at least
\(m/159\); if two were smaller, their pair sum would be too small.
Their supports sum to at least \(1-2^{-m/318}\), by TD08,
and this is at least \(15/16\).
The last nonidentity projected entry retains the universal
Heisenberg quarter floor, including if its label is identity.
Thus
\[
 2+\frac6e
   >\frac{22}{25}+\frac{15}{16}+\frac14
    =\frac{827}{400},
       \qquad e<\frac{800}{9}.
 \tag{HFS.12}
\]
This contradicts \(e\ge2^m\ge2^{1272}>800/9\).
There is therefore no \(D_i\ge160\) in this range.
When \(m<1272\), every \(D_i\le2m<2544\).
Together with \(D_i=1\) at omitted entries, this proves the
four-entry assertion in its stated stronger standalone scope.

\emph{Step 4: every remaining prime and low-exponent branch.}
The actual noncyclic projective image contains \(Q\), so the
generating product-one projected tuple has \(k\ge3\).
For \(p=2\), the accepted TD59 theorem says that \(m>80\)
forces \(k=3\) or4.  These cases have just been treated.
If \(m\le80\), then \(D_i\le160\le C(M)\).

For \(p=3\), the accepted TD57 alternatives give either
\[
      k\in\{3,4\},\qquad
      m\le4,\qquad\text{or}\qquad
                   3^m\le\frac{20M}{7}.
 \tag{HFS.13}
\]
The first two cases are covered by Steps 2--3 and \(D_i\le8\).
In the last case,
\[
        2^m\le3^m\le20M/7<4M\le2^{2M}.
 \tag{HFS.14}
\]
Here \(4M\le4^M\) for positive integers \(M\): equality holds
at \(M=1\), and the inequality propagates because
\(4(M+1)\le16M\).
It follows that \(m<2M\), so \(D_i<4M\le64M^2\).

For every \(p\ge5\), TD57 gives either \(k=3,4\) or \(m\le3\).
The latter case has \(D_i\le6\).
This proves \(D_i\le C(M)\) in every case, without dropping
any low-exponent alternative.
All uses of TD57 and TD59 retain primitivity of the actual
projective factor \(P\).  No source budget is applied to \(W\)
or to \(U\) as a separate sphere coefficient.

\emph{Step 5: label orders and the exact projective kernel.}
TD55, with the same actual \(M\)-floor and original marks, gives
\(p\le B(M)\).
For any invertible \(g\in\operatorname{GL}(Q)\) of minimal-polynomial
degree \(D\), the cyclic subspace constructed in TD61 has dimension
\(D\) and the same minimal polynomial as \(g\).
The cyclic group generated by \(g\) acts faithfully on that subspace:
if \(g^j\) is identity there, its polynomial \(X^j-1\) is divisible
by the common minimal polynomial and hence kills all of \(Q\).
Thus, including the defining-characteristic part of the order,
\[
  \operatorname{ord}(g)\le|\operatorname{GL}_D(\mathbf F_p)|
                  <p^{D^2}.
 \tag{HFS.15}
\]
No semisimplicity of \(g\) is required.

The exact projective centralizer statement TD05 identifies the
kernel of \(P\to\operatorname{Sp}(Q)\) with \(Q\) itself.
Since \(Q\) has exponent \(p\), for every original index one has
\[
 \operatorname{ord}(y_i)\mid p\,\operatorname{ord}(g_i),
 \qquad
 \operatorname{ord}(y_i)
       <p^{D_i^2+1}\le B(M)^{C(M)^2+1}=O(M).
 \tag{HFS.16}
\]
This includes \(g_i=I\), nonidentity Weyl translations and
projectively identity entries.  It is a statement about the
actual projective quotient, not the order of a chosen scalar lift.

Finally \(e\ge a\ge2\), so the whole-original floor gives
\[
             \frac rM\le2+\frac6e\le5,
                         \qquad r\le5M.
 \tag{HFS.17}
\]
The finite set of ordered lists of length at most \(5M\)
with integer entries between1 and \(O(M)\) contains every
selected projective signature.  Counting all these lists, including
many which fail every geometric or generation condition, already
gives a finite enumeration.  No realization assertion is used.
\end{proof}

\paragraph{Actual minimal-normal extraction and the cutoff application.}
In the complete TD12 recursion, each Heisenberg coarse quotient
comes from a chosen minimal normal elementary abelian subgroup.
Its image \(Q\) is faithful.  A nonzero proper \(H\)-invariant
\(\mathbf F_p\)-subspace of \(Q\) would pull back to a nonzero
proper normal subgroup of that chosen minimal normal subgroup,
contradicting minimality.  Hence its label module is irreducible.
The same recursion proves that the coarse projective factor is
primitive and that its action is an actual factorwise quotient
of the original coefficient, with the entire complement retained.
These facts supply precisely the extra hypotheses used here.

If a separate complete alternating cutoff supplies the whole
original floor \(M_0=\max\{80,N-2\}\), then the present explicit
integer bounds apply simultaneously to every such Heisenberg
coarse extraction of the same source.  In particular \(N=N_*\)
may be used in the exact minimal, naturally nonsplit,
imaginary-quadratic, rank-nine comparison class of accepted TD54.
No such cutoff is asserted outside that class.

\paragraph{Remaining boundary.}
The finite-signature argument supplies no bound on the order of
the generated projective quotient, the dimension \(2m\) of its
irreducible label module, or the number of possible representations
or Nielsen classes.  It does not assert that any of these quantities
is actually unbounded among geometric sources.
Original scalar periods are still unrestricted by this theorem.
No group, cover, original-field realization, prescribed polarization,
period-map image, or effective detecting cycle is constructed.
